% =====================================================================
%  2 Testprotokoll  (Aufgabenstellung: „Für jeden Testfall ist Folgendes
%  zu protokollieren: eindeutige Testfall-ID, Kurzbeschreibung,
%  Vorbedingungen, Eingabedaten und/oder auszuführende Aktionen,
%  erwartetes Ergebnis, tatsächliches Ergebnis.“)
%  Je Testfall ein Block mit dem Makro \testfall (alle Pflichtfelder).
% =====================================================================
\newcommand{\testfall}[8]{%
  % #1 ID  #2 Kurzbeschreibung  #3 Bezug (US/NFA/GR)  #4 Teststufe
  % #5 Vorbedingungen  #6 Eingabedaten/Aktionen  #7 erwartet  #8 tatsächlich
  \begin{longtable}{@{}L{4.0cm}L{11.2cm}@{}}
    \toprule
    \textbf{Testfall-ID} & \textbf{#1} \\
    \midrule
    Kurzbeschreibung & #2 \\
    Bezug & #3 \\
    Teststufe & #4 \\
    Vorbedingungen & #5 \\
    Eingabedaten/Aktionen & #6 \\
    Erwartetes Ergebnis & #7 \\
    Tatsächliches Ergebnis & #8 \\
    \bottomrule
  \end{longtable}}

\section{Testprotokoll}
\TODO{Kurze Hinführung; Konventionen der IDs (z. B. TC-U-01 Unit, TC-I-01 Integration, TC-S-01 System)}

\subsection{Unit-Tests}
\testfall{TC-U-01}{\TODO{}}{\TODO{US01}}{Unit-Test}{\TODO{}}{\TODO{}}{\TODO{}}{\TODO{bestanden/nicht bestanden + Beobachtung}}

\subsection{Integrationstests}
\testfall{TC-I-01}{\TODO{}}{\TODO{}}{Integrationstest}{\TODO{}}{\TODO{}}{\TODO{}}{\TODO{}}

\subsection{Systemtests}
\testfall{TC-S-01}{\TODO{}}{\TODO{}}{Systemtest}{\TODO{}}{\TODO{}}{\TODO{}}{\TODO{}}
